% Chapter 7 — Conclusion
% Converted from thesis_notes/draft_ch7_conclusion.md
\chapter{Conclusion}
\label{chap:conclusion}
\lhead{\emph{Conclusion}}

\section{Summary of work}
\label{sec:concl-summary}

This thesis identified a compact, physically interpretable model of a DC gearmotor driven by an L298N H-bridge, measured by a quadrature encoder on an Arduino Uno R4 Minima. The model is a \textbf{cascade}---PWM command\,$\to$\,static driver-voltage block\,$\to$\,motor dynamic block\,$\to$\,RPM (Model~C)---and each block was identified through a sequence of bench and analysis phases: a static sweep and deadzone characterisation (Phase~1.1/1.2), an automated step-response campaign (Phase~2.1), a leave-one-out validation of that campaign (Phase~2.2, analysis only), and a gap-fill bench session two weeks later that added operating points and bounded session-to-session drift.

\section{Principal findings}
\label{sec:concl-findings}

\textbf{The cascade decomposition is empirically justified, not assumed.} The plant's large forward/reverse asymmetry localises almost entirely to the driver: the PWM\,$\to$\,voltage map differs by direction by a factor of up to 2.2 near the deadzone edge, while the motor's voltage\,$\to$\,speed gain is direction-symmetric to within 9\,\% ($K_{\mathrm{fwd}} = 26.94$, $K_{\mathrm{rev}} = 24.52$~rpm/V). Independently, the only measurable session drift sits on the driver side as well, with the motor block reproducible to ${\approx}1\,\%$. Asymmetry and instability therefore both belong to one identifiable block, which is precisely what makes the cascade the right structure (\cref{chap:static}).

\textbf{The dynamics are region-dependent---Model~C, not a single LTI system.} Acceleration is well described by FOPDT whose time constant follows a broad bathtub in PWM (${\approx}81$\,ms at the mid-range minimum, rising to ${\approx}200$\,ms near breakaway and ${\approx}150$\,ms near full drive). Deceleration is a separate regime: its time constant is roughly twice the acceleration value and is non-monotone in the target, peaking at ${\approx}321$--$366$\,ms for stops that coast through the high-friction deadzone. The forward/reverse asymmetry even reverses sign between acceleration and deceleration---a small but real departure from a perfectly clean cascade. A single global time constant is misspecified by ${\approx}2\times$ and is rejected (\cref{chap:dynamic}).

\textbf{The model validates within its grid, and validation surfaced a cold-start effect.} Leave-one-out cross-validation---adopted because no held-out dataset ever existed---gives a predictive excess over the in-sample fit of $+1.06$\,rpm for acceleration and $+0.10$\,rpm for deceleration on the Phase~2.1 grid. The $10\times$ gap traced to cold-start stochasticity: acceleration from rest catches the motor cold, with stochastic friction breakaway, whereas deceleration always begins from a warm, settled state. The gap-fill session confirmed this falsifiable prediction directly---with a warm motor, acceleration excess collapsed to $-0.05$\,rpm (\cref{chap:validation}).

\textbf{Session discipline is a methodological result in its own right.} Re-running the Phase~2.1 conditions two weeks later revealed structured drift---a cold-start warm-up transient, a persistent low-speed deceleration slowdown, and a persistent ${\approx}20\,\%$ reverse high-PWM speed-up. The practical consequence is firm: calibration does not transfer across sessions for deceleration, so any later work must recalibrate in situ rather than reuse stored parameters.

\section{Limitations}
\label{sec:concl-limits}

The static map rests on a single sweep (no run02/run03), supported by a two-point spot-check and the implicit cross-validation of the step-response endpoints rather than by full repeat sweeps. The dynamic tables are a two-session composite, honestly labelled, because the bathtub and target-dependence pictures require both sessions. The FOPDT fits are \emph{effective} at the operating-point extremes and for coast-to-rest decel, where the true behaviour is friction-dominated and not first-order, and the fitted dead time is consequently an effective parameter rather than a physical latency. Validation is interpolation within the calibration grid only; it tests neither extrapolation nor---given the documented drift---uncorrected cross-session prediction. One condition, reverse $240\to100$, is irreducibly noisy through the deadzone.

\section{Future work}
\label{sec:concl-future}

The thesis defers its culminating step---closed-loop validation---for which \cref{chap:closedloop} gives a concrete plan: a multi-transition simulator extracted into a reusable module, a PID firmware variant with deadzone feedforward, a global-LTI baseline to beat, and a self-contained session that calibrates in situ before comparing cascade-aware against baseline control. Independent open-loop refinements remain available: a friction-explicit ``Model~D'' (Coulomb/Stribeck) for the deadzone and coast regimes, a higher-order dynamic block near full drive, extension of the operating grid (below breakaway, to PWM~=~255, and to other deceleration start points), an instrumented study of the unexplained reverse-$\tau$ deflation, and routine ambient/driver temperature logging.

\section{Closing}
\label{sec:concl-closing}

The open-loop system identification is complete and validated within its calibration grid. The contribution is a cascade model whose structure is earned from the data---asymmetry and drift isolated to the driver, region- and operating-point-dependent dynamics in the motor---together with a validation that not only quantified predictive error but exposed and then confirmed a cold-start mechanism, and a clear methodological lesson about in-session calibration. What remains is to close the loop and show the model earns its complexity in control; the groundwork for that demonstration is now fully in place.

\chaptersources{Chapters~\ref{chap:hardware}--\ref{chap:closedloop} and their underlying processed CSVs \textperiodcentered{} status-doc \S4 (final model state), \S5 (limitations), \S6 (future work) \textperiodcentered{} \texttt{thesis\_notes/log.md} (full session narrative).}
